\begin{figure*}[t]
\textbf{Paper-level field} (\texttt{bibliographic.title}):
\begin{lstlisting}[style=codeblock]
{
  "title": {
    "description": "Main title of the scientific document itself.",
    "type": "string|null",
    "cardinality": {
      "min": 0,
      "max": 1
    },
    "instructions": [
      "Copy the document title verbatim as printed on the first page or running header. Preserve casing and punctuation."
    ]
  }
}
\end{lstlisting}
\textbf{Entity with sub-schema} (\texttt{samples}, abridged):
\begin{lstlisting}[style=codeblock]
{
  "samples": {
    "description": "Samples studied through experimental measurements reported or reanalyzed in the document. A sample entry can represent a single specimen or a measured group when only the group is identified by the document.",
    "type": "list[sample]",
    "cardinality": {
      "min": 0,
      "max": 50
    },
    "instructions": [
      "Follow this decision routine to determine sample entries.",
      "..."
    ],
    "item_schema": {
      "id": {
        "type": "string",
        "description": "Stable local sample identifier.",
        "instructions": [
          "Assign IDs in sequence: sample_1, sample_2, sample_3, etc.",
          "..."
        ]
      },
      "sample_name": {
        "type": "string",
        "description": "Name, label, or concise identifying phrase for the sample.",
        "instructions": [
          "Follow this decision routine to extract sample_name.",
          "..."
        ]
      },
      "...": "5 further subfields: sample_description, chemical_composition, sample_state, sample_mass, supplier_or_creator"
    }
  }
}
\end{lstlisting}
  \caption{Two excerpts from the released two-layer schema: a paper-level
  field shown in full, and an entity collection with
  its \texttt{item\_schema} of subfields. Each field carries a
  \texttt{description} and \texttt{type}; top-level fields also declare
  a \texttt{cardinality}, and most fields carry an \texttt{instructions}
  array (removed by the no-instruction ablation of
  Section~\ref{sec:instruction-ablation}). The \texttt{samples} entry is
  abridged: long instruction routines and the remaining subfields are
  elided (marked \texttt{...}), and the complete schema with all
  per-field definitions ships with the benchmark.}
  \label{fig:schema-excerpt}
\end{figure*}
